\section{מגבלות}

גם לאחר בדיקות היציבות וההרחבות שבוצעו, יש מספר מגבלות שחשוב לשמור מולן על ניסוח זהיר.

\begin{itemize}
    \item \textbf{מספר קטן של תכונות.} ארבע התכונות שנבחרו מתארות בעיקר כיוון, תנודתיות, טווח מסחר ונפח. הן מתאימות לניתוח מצב השוק, אך ייתכן שאינן מכילות מספיק מידע לחיזוי כיוון של יום בודד. הוספת נתונים מאקרו־כלכליים, מידע מאופציות או \tech{order flow} עשויה לשנות את התוצאה, אך גם מגדילה את הסיכון ל־\tech{overfitting} ולחיפוש יתר אחר תוצאה טובה.

    \item \textbf{בדיקת זמן מלאה רק על \asset{SPY}.} לכל נכס קיימת תקופת \tech{Test} נפרדת, אך בדיקת \tech{walk-forward} בשלושה חלונות בוצעה רק על SPY. לכן העמידות בזמן נבדקה לעומק עבור נכס הייחוס ולא עבור כל תשעת הנכסים.

    \item \textbf{בחירת מספר המצבים.} בחירה לפי \tech{Validation log-likelihood} העדיפה ברוב הנכסים \(K=4\). גם לאחר בדיקת יציבות בין אתחולים, ייתכן שמספר המצבים המתאים משתנה לאורך זמן או בין סוגי נכסים.

    \item \textbf{שמות המצבים הם פרשנות.} שמות כמו \tech{Calm} או \tech{Crisis-like} ניתנו לאחר בחינת הסטטיסטיקות. ה־HMM עצמו אינו יודע מהו ``משבר'' ואינו מקבל תוויות כלכליות בזמן האימון.

    \item \textbf{האימות מול \asset{VIX} חלקי.} המצב הנדיר 2 של \asset{SPY} לא הופיע בתקופת ה־\tech{Test} שבה בוצעה ההשוואה ל־VIX. לכן לא ניתן לבדוק את הפרשנות שלו באמצעות המדד החיצוני באותו חלון.

    \item \textbf{מודלי ה־\tech{Supervised Learning} אינם תחרות מכוונת־ביצועים.} שלושת המסווגים נבחרו כמודלי ייחוס שקופים עם \tech{hyperparameters} קבועים. לא בוצע חיפוש נרחב אחר התצורה הטובה ביותר לכל מסווג.

    \item \textbf{כמות נתוני האימון אינה זהה לחלוטין.} לאחר בחירת \(K\) וה־\tech{seed}, ה־HMM הסופי וכללי הייחוס הפשוטים משתמשים בכל היסטוריית \tech{Train+Validation} שהייתה זמינה לפני ה־\tech{Test}. מודלי ה־\tech{Supervised Learning} אומנו על \tech{Train} בלבד. לכן הדירוג ביניהם הוא אינדיקציה שימושית, אך אינו ניסוי שבו לכל המודלים ניתנה בדיוק אותה כמות מידע.

    \item \textbf{אין בדיקת רווחיות מסחר.} DPA, MAE או RMSE אינם שקולים לתשואה כספית. לא בוצע \tech{backtest} הכולל עמלות, החלקת מחיר, מיסוי או גודל פוזיציה. לכן אין להסיק מהתוצאות המלצת מסחר או רווחיות צפויה.
\end{itemize}

\section{מסקנות}

הפרויקט התחיל משאלה פשוטה יחסית: האם \tech{Gaussian HMM} יכול לזהות מצבים חבויים בשוק ולנצל אותם כדי לשפר את חיזוי כיוון התשואה ביום הבא. לאחר הרחבת הניסוי לתשעה נכסים, כמה ערכי \(K\), מספר אתחולים, מודלי ייחוס ובדיקות יציבות, התקבלה תשובה שמפרידה בין חיזוי לבין ייצוג מצב השוק.

מבחינת חיזוי, התוצאה חלשה יחסית. כלל הייחוס הפשוט המבוסס על ממוצע העבר השיג DPA ממוצע של \pct{54.54}, לעומת \pct{53.06} לגרסה הרכה של HMM ו־\pct{52.06} לגרסה הקשיחה. גם מודלי \tech{Supervised Learning} נשארו באותו סדר גודל, ובדיקת ה־\tech{walk-forward} על \asset{SPY} הראתה פערים קטנים ולא עקביים בין התקופות. גם במדדי MAE ו־RMSE על SPY ה־HMM כמעט זהה לכלל ה־\tech{Naive}. לכן אין בסיס לטעון שה־HMM מספק יתרון יציב בחיזוי היום הבא.

מבחינת מבנה השוק, התוצאה מעניינת יותר. ב־\asset{SPY} התקבלו ארבעה מצבים בעלי מאפיינים שונים של תשואה, תנודתיות, טווח יומי, \tech{drawdown} ומשך שהייה. מצב נדיר אחד הציג מאפייני לחץ ברורים. חלוקת המצבים הייתה יציבה מאוד בין אתחולים, ו־\asset{VIX} חיצוני קיבל רמות שונות בחלק מהמצבים. בנוסף, הקורלציה בין SPY לנכסים אחרים השתנתה בהתאם למצב שזוהה.

לכן המסקנה המרכזית היא שה־HMM מתאים בפרויקט זה יותר ל־\textbf{תיאור ואמידה של מבנה שוק מותנה} מאשר ל־\textbf{חיזוי נקודתי}. הוא אינו אומר בביטחון מה תהיה התשואה מחר, אך הוא כן מספק הקשר הסתברותי: האם השוק נמצא בסביבה רגועה או תנודתית, כמה זמן מצב כזה נוטה להימשך, ועד כמה המודל בטוח בזיהוי.

הרחבה אפשרית היא להשתמש בווקטור ה־\tech{posterior} של ה־HMM כחלק ממערכת \tech{Regime-aware Reinforcement Learning}, ולבדוק האם מידע על מצב השוק מסייע בקבלת החלטות לאורך זמן תחת עלויות עסקה ובדיקת \tech{walk-forward}. זהו ניסוי חדש ונפרד, ולא מסקנה שניתן להסיק ישירות מהתוצאות הקיימות.
